\section{Discussion and Future Work}
\label{sec:discussion}

\paragraph{Where SignedKAN sits.}
On signed-graph link-sign prediction, SignedKAN's
trained-from-scratch numbers are competitive with engineered
baselines: AUC $\approx\!0.85$--$0.88$ on Bitcoin Alpha and OTC
(comparable to SiNE, DSHGNN, and the lower end of SDGNN's reported
range), $\sim\!0.05$--$0.07$ below SGCN's published number.
Macro-$F_1$ on Bitcoin OTC reaches $0.81$ under the bilinear
endpoint head extension. None of these comparisons is under matched
training/split conditions, so they are informative rather than
definitive; the paper reports the regime under which it actually ran.

\paragraph{The differentiating contribution.}
The four-part claim of Section~\ref{sec:pruning} is what
distinguishes SignedKAN from engineered signed-GNN baselines:
(i) competitive trained-from-scratch AUC; (ii) 77--78\% pruning
resilience with AUC \emph{improving} ($+0.016$ to $+0.021$);
(iii) 80--91\% sinusoidal symbolic distillation, basis-invariant
across three spline parameterisations; (iv) sub-minute training time
on a single GPU. Engineered baselines (SGCN, SDGNN) cannot be
distilled this cleanly and offer no comparable compression
guarantee. The model-compression and interpretability properties
we measure are not architectural ornaments — they are the
load-bearing claim.

\paragraph{Limitations.}
We do not provide a formal expressivity result for the three-branch
sign-conditioned aggregation, nor a closed-form analysis of the
training dynamics that drive activations toward sinusoidal forms.
The cross-paper AUC gap to SGCN is structural to the experimental
regime (no pre-training, no auxiliary supervision) — closing it
under matched conditions is left to a journal extension that can
afford a re-implementation of the engineered baselines under our
split protocol. The Slashdot benchmark is bottlenecked by GPU memory
at $h\!=\!16+$ and was reported here only at $h\!=\!8$.

\paragraph{Future work — clique clusters and higher-order balance.}
The signed hypergraph used here is built from $3$-uniform triads
under Cartwright--Harary balance. The natural next axis is
\emph{$k$-clique clusters}: signed cliques of size $k\!\geq\!3$
generalise the balance criterion (Davis~1967, weakly-balanced
$k$-cycles) and extend the hypergraph IR's arity. Because the
SignedKAN layer is already arity-agnostic over the sub-aggregation
loop, supporting $k$-uniform hyperedges only requires extending the
construction step. A complementary direction is a \emph{clique-level
entropy regulariser}: per-clique participation entropy as an
additional term in the spectral-entropy schedule of
Section~\ref{sec:experiments}, providing structural prior at the
hyperedge-cluster level rather than only the
node-embedding-spectrum level.

\paragraph{Hypergraph-IR integration.}
The SignedKAN architecture maps to a Tier-2 emission target in a
canonical-hypergraph IR back end with a torch-dataflow code path.
Defining \texttt{lyr.signedkan\_layer} in the meta library and
adding a \texttt{SignedKANBlock} to the block-helper library would
let downstream applications declare SignedKAN topologies in IR
source and emit the corresponding PyTorch module automatically.
We sketch this integration but defer the implementation.